\scanpage{193}

Si tous les $X_{ji}$ sont nuls, alors ($V$ étant indécomposable) il n'y a qu'un seul $V_n$ non nul, qui est de dimension un. On désignera la repr. ainsi obtenue (définie par la repr. $\sigma_n$ de $H=C^*$ en composant avec $G\longrightarrow H$) par \uline{$\sigma_n$}. S'il existe un $X_{\beta\alpha}$ non nul, alors \inserted{$\beta-\alpha=p$, donc} on a $V=V'+V^{\unsure{\prime\prime\prime|\prime\prime}}$,

où $V'=\sum_{\substack{n\equiv\alpha\\(p)}}V_n$, $V''=\sum_{\substack{n\not\equiv\alpha\\(p)}}V_n$

sont stables sous $\unsure{\mathfrak g|\mathscr G}$, et $V'\ne0$, donc $V'=V$. \deleted{Soit} On a donc, si \deleted{\unsure{$p$ est le plus}} \inserted{$V_n,V_{n'}\ne0$ : $n'-n\equiv0\ (p)$.}

\inserted{Soit $n$ le plus} petit entier tel que $V_{\unsure{n'|n}}\ne0$, alors les $V_m$ éventuellement non nuls \deleted{\unsure{du}} sont de la forme $V_{n+pk}$ $k=0,1,\ldots$. Soit $k$ le plus grand entier tel que $V_{n+pk}\ne0$, \deleted{\unsure{alors}} on a $k\geq1$. Soit $u_i=X_{n+p(i+1),n+pi}$,

\deleted{\unsure{on a} \unsure{$a=0$} $V$} $E_i=V_{n+pi}$, on a donc une suite d'applications

\diagram{\begin{tikzcd}[ampersand replacement=\&,column sep=small] E_0\arrow[r,"u_0"] \& E_1\arrow[r,"u_1"] \& E_2 \& \ldots\arrow[r,"u_{k-1}"] \& E_k. \end{tikzcd}}

\uline{Lemme 1}. Les $u_i$ sont \uline{injectifs}. \inserted{$(0\leq i\leq k-1)$} On va le prouver par récurrence sur $i$, c'est trivial pour $i=-1$ : $0=E_{-1}\longrightarrow E_0$. Supposons

% archive sheet number: 192
